\documentclass[11pt]{article}
\usepackage[a4paper,margin=28mm]{geometry}
\usepackage{amsmath,amssymb,amsthm}
\newtheorem{theorem}{Theorem}
\newtheorem{lemma}[theorem]{Lemma}
\newtheorem{corollary}[theorem]{Corollary}
\title{A four-face die cannot belong to five permutation-fair dice}
\author{Research note, 30 September 2026}
\date{}
\begin{document}
\maketitle
\noindent\textbf{Status.} Complete elementary proof, independently reviewed.
The symbolic determinant identity is checked by the accompanying exact
script. Numerical exploration suggested the proof but is not used in it.

\begin{theorem}\label{thm:no4}
No permutation-fair collection of five dice contains a die with exactly
four faces. Consequently, in conjunction with the universal $n-1$ bound,
every die in such a collection has at least five faces.
\end{theorem}

We prove a stronger local statement using comparison moments only through
degree three. No fourth-order comparison moment or global order condition
is needed.

\begin{lemma}\label{lem:local}
There do not exist four nondecreasing rational-valued functions
$p_j:\{1,2,3,4\}\to[0,1]$, $j\in[4]$, such that
\[
 \frac14\sum_{q=1}^4\prod_{j\in S}p_j(q)=\frac1{|S|+1}
 \qquad(S\subseteq[4],\ |S|\leq3).
\]
\end{lemma}
\begin{proof}
Put $y_j=2p_j-1$. The assumed moments imply
\[
 \mathbb E y_j=0,\qquad
 \mathbb E(y_i y_j)=\frac13\ (i\ne j),\qquad
 \mathbb E(y_i y_j y_k)=0\quad(i,j,k\text{ distinct}),
\]
where expectation is uniform on the four ordered values of $q$.
Use the orthonormal basis of the mean-zero subspace consisting of the
columns of
\[
 H=\begin{pmatrix}
 -1&-1&-1\\
 -1& 1& 1\\
  1&-1& 1\\
  1& 1&-1
 \end{pmatrix},
 \qquad \frac14H^TH=I_3.
\]
Write $y_j=H(A_j,B_j,C_j)^T$. All coefficients are rational.
The three adjacent differences of $y_j$ are
\[
 2(B_j+C_j),\qquad2(A_j-B_j),\qquad2(B_j-C_j).
\]
Monotonicity is therefore equivalent to
\begin{equation}\label{eq:cone}
 A_j\geq B_j\geq|C_j|.
\end{equation}
Moreover $A_j>0$, since $A_j=0$ would make $y_j=0$, contradicting
its nonzero pairwise inner products. Orthonormality gives
\begin{equation}\label{eq:pair}
 A_iA_j+B_iB_j+C_iC_j=\frac13\qquad(i\ne j).
\end{equation}

Put $b_j=B_j/A_j\geq0$ and $c_j=C_j/A_j$.
The product of the three distinct columns of $H$ is identically $-1$;
all other third-order products of its columns have expectation zero.
Consequently the triple-moment condition, divided by the positive product
$A_iA_jA_k$, is
\[
 c_i(b_j+b_k)+c_j(b_i+b_k)+c_k(b_i+b_j)=0
 \qquad(i,j,k\text{ distinct}).
\]
Collect the four equations according to their omitted index. They say
$M(b)c=0$, where
\[
 M(b)_{ii}=0,\qquad
 M(b)_{ij}=\sum_{k\notin\{i,j\}}b_k\quad(i\ne j).
\]
Writing $e_r$ for elementary symmetric functions, direct expansion gives
\begin{equation}\label{eq:det}
 \det M(b)=-4\bigl(e_1(b)e_3(b)-4e_4(b)\bigr)
 =-4\sum_{\substack{i,j,k\text{ distinct}\\j<k}}b_i^2b_jb_k.
\end{equation}
The sum on the right has twelve terms. If at least three of the $b_j$
are positive, it is strictly positive, so $M(b)$ is invertible. Hence
all $C_j=0$.

In this case the four rational vectors $v_j=(A_j,B_j)\in\mathbb Q^2$
have pairwise inner products $1/3$. Their Gram matrix is
\[
 G=D+\frac13\boldsymbol1\boldsymbol1^T,
 \qquad D_{jj}=A_j^2+B_j^2-\frac13.
\]
Since $\operatorname{rank}G\leq2$, the diagonal matrix $D$ has rank at
most three. Thus some $A_j^2+B_j^2=1/3$. This is impossible over
$\mathbb Q$: clearing denominators in $a^2+b^2=1/3$ gives
$3(u^2+v^2)=w^2$ for coprime integers $u,v,w$. It first forces $3\mid w$,
then $u^2+v^2\equiv0\pmod3$, so $3\mid u,v$, a contradiction.

It remains to consider at most two positive $b_j$. At least two indices,
say $i$ and $j$, have $B_i=B_j=0$, and \eqref{eq:cone} also gives
$C_i=C_j=0$. Choose a third index $k$. Equation~\eqref{eq:pair} gives
\[
 A_iA_j=A_iA_k=A_jA_k=\frac13.
\]
Positivity implies $A_i=A_j=A_k$ and $A_i^2=1/3$, again impossible for
rational $A_i$. Both cases contradict the assumed functions.
\end{proof}

\begin{proof}[Proof of Theorem~\ref{thm:no4}]
Suppose that the distinguished die $D_0$ has four faces, in increasing
order. For the other four dice put $p_j(q)=\Pr(D_j<D_0(q))$.
These functions are rational and nondecreasing. Permutation fairness,
restricted to $D_0$ and any set $S$ of the other dice, gives
\[
 \frac14\sum_q\prod_{j\in S}p_j(q)
  =\Pr(D_j<D_0\text{ for every }j\in S)=\frac1{|S|+1}.
\]
Lemma~\ref{lem:local} is a contradiction.
\end{proof}

\paragraph{Consequence for the exceptional-die optimum.}
For $g(n)$, the minimum number of faces of a prescribed die with all other
dice unrestricted, the independently established results now read
\[
 g(2)=1,\qquad g(3)=2,\qquad g(4)=3,\qquad g(5)\geq5,
 \qquad n-1\leq g(n)\leq Cn^2.
\]
In particular, the tempting equality $g(n)=n-1$ for every $n$ is false.
The exact value of $g(5)$ and the asymptotic order remain open here.
\end{document}
